% \tikzset{
%   >=latex, % for default LaTeX arrow head
%   node/.style={thick,circle,draw=myblue,minimum size=22,inner sep=0.5,outer sep=0.6},
%   node in/.style={node,green!20!black,draw=mygreen!30!black,fill=mygreen!25},
%   node hidden/.style={node,blue!20!black,draw=myblue!30!black,fill=myblue!20},
%   node convol/.style={node,orange!20!black,draw=myorange!30!black,fill=myorange!20},
%   node out/.style={node,red!20!black,draw=myred!30!black,fill=myred!20},
%   connect/.style={thick,mydarkblue}, %,line cap=round
%   connect arrow/.style={-{Latex[length=4,width=3.5]},thick,mydarkblue,shorten <=0.5,shorten >=1},
%   node 1/.style={node in}, % node styles, numbered for easy mapping with \nstyle
%   node 2/.style={node hidden},
%   node 3/.style={node out}
% }
\tikzset{
  >=latex, % for default LaTeX arrow head
  node/.style={thick,circle,draw=black,minimum size=22,inner sep=0.5,outer sep=0.6},
  node in/.style={node,draw=black,fill=myblue!20}, % Input nodes (gray)
  node hidden/.style={node,draw=black,fill=myblue!20}, % Hidden nodes (darker gray)
  node out/.style={node,draw=black,fill=myblue!20}, % Output nodes (darkest gray)
  connect/.style={thick,black}, %,line cap=round
  connect arrow/.style={-{Latex[length=4,width=3.5]},thick,black,shorten <=0.5,shorten >=1},
  node 1/.style={node in}, % node styles, numbered for easy mapping with \nstyle
  node 2/.style={node hidden},
  node 3/.style={node out}
}
\def\nstyle{int(\lay<\Nnodlen?min(2,\lay):3)} % map layer number onto 1, 2, or 3

\begin{tikzpicture}[x=2.2cm,y=1.4cm]
  \message{^^JNeural network, shifted}
  \readlist\Nnod{4,5,5,5,4} % array of number of nodes per layer
  \readlist\Nstr{31,128,128,128,4} % array of string number of nodes per layer
  \readlist\Cstr{\strut o,h^{(\prev)},h^{(\prev)},h^{(\prev)},a} % array of coefficient symbol per layer
  \def\yshift{0.5} % shift last node for dots
  
  \message{^^J  Layer}
  \foreachitem \N \in \Nnod{ % loop over layers
    \def\lay{\Ncnt} % alias of index of current layer
    \pgfmathsetmacro\prev{int(\Ncnt-1)} % number of previous layer
    \message{\lay,}
    \foreach \i [evaluate={\c=int(\i==\N); \y=\N/2-\i-\c*\yshift;
                 \index=(\i<\N?int(\i):"\Nstr[\lay]");
                 \x=\lay; \n=\nstyle;}] in {1,...,\N}{ % loop over nodes
      % NODES
      \node[node \n] (N\lay-\i) at (\x,\y) {$\Cstr[\lay]_{\index}$};
      
      % CONNECTIONS
      \ifnum\lay>1 % connect to previous layer
        \foreach \j in {1,...,\Nnod[\prev]}{ % loop over nodes in previous layer
          \draw[connect,white,line width=1.2] (N\prev-\j) -- (N\lay-\i);
          \draw[connect] (N\prev-\j) -- (N\lay-\i);
          %\draw[connect] (N\prev-\j.0) -- (N\lay-\i.180); % connect to left
        }
      \fi % else: nothing to connect first layer
      
    }
    % Only show vertical dots for layers that are not the last layer
    \ifnum\lay<\Nnodlen
      \path (N\lay-\N) --++ (0,1+\yshift) node[midway,scale=1.5] {$\vdots$};
    \fi
  }
  
  % LABELS
  \node[above=1,align=center] at (N1-2.90) {input\\[-0.2em]layer};
  \node[above=1,align=center] at (N3-2.90) {hidden layers};
  \node[above=2,align=center] at (N3-2.90) {\LARGE $\pi$-Network};
  \node[above=1,align=center] at (N\Nnodlen-2.90) {output\\[-0.2em]layer};
\end{tikzpicture}

% \tikzset{
%   >=latex, % for default LaTeX arrow head
%   node/.style={thick,circle,draw=black,minimum size=22,inner sep=0.5,outer sep=0.6},
%   node in/.style={node,draw=black,fill=gray!20}, % Input nodes (gray)
%   node hidden/.style={node,draw=black,fill=gray!40}, % Hidden nodes (darker gray)
%   node out/.style={node,draw=black,fill=gray!60}, % Output nodes (darkest gray)
%   connect/.style={thick,black}, %,line cap=round
%   connect arrow/.style={-{Latex[length=4,width=3.5]},thick,black,shorten <=0.5,shorten >=1},
%   node 1/.style={node in}, % node styles, numbered for easy mapping with \nstyle
%   node 2/.style={node hidden},
%   node 3/.style={node out}
% }
% \def\nstyle{int(\lay<\Nnodlen?min(2,\lay):3)} % map layer number onto 1, 2, or 3

% \begin{tikzpicture}[x=2.2cm,y=1.4cm]
%   \message{^^JNeural network, shifted}
%   % ** EDITED: Nodes per layer: 9 Input, 128 (dots), 128 (dots), 128 (dots), 4 Output **
%   % We will represent the three hidden layers as three columns of 5 nodes each, 
%   % and use the labels to imply the full size (128).
%   \readlist\Nnod{9,9,9,9,4} 
%   \readlist\Nstr{9,128,128,128,4} % String representation for labels
  
%   % ** EDITED: Node labels for clarity **
%   \def\yshift{0.5} % shift last node for dots
%   \readlist\Cstr{o,h,h,h,a} % Variable name: observation (o), hidden (h), action (a)
  
%   \message{^^J 	Layer}
%   \foreachitem \N \in \Nnod{ % loop over layers
%     \def\lay{\Ncnt} % alias of index of current layer
%     \pgfmathsetmacro\prev{int(\Ncnt-1)} % number of previous layer
%     \message{\lay,}
%     \foreach \i [evaluate={\c=int(\i==\N); \y=\N/2-\i-\c*\yshift;
%                  \index=(\i<\N?int(\i):\Nstr[\lay]);
%                  \x=\lay; \n=\nstyle;}] in {1,...,\N}{ % loop over nodes
%       % NODES
%       \node[node \n] (N\lay-\i) at (\x,\y) {$\Cstr[\lay]_{\index}$};
      
%       % CONNECTIONS
%       \ifnum\lay>1 % connect to previous layer
%         \foreach \j in {1,...,\Nnod[\prev]}{ % loop over nodes in previous layer
%           \draw[connect,line width=1.2] (N\prev-\j) -- (N\lay-\i); % Removed white background for connectors
%           \draw[connect] (N\prev-\j) -- (N\lay-\i);
%         }
%       \fi % else: nothing to connect first layer
      
%     }
%     % DOTS: Only for the hidden layers (2, 3, 4)
%     \ifnum\lay>1
%       \ifnum\lay<\Nnodlen
%         \path (N\lay-\N) --++ (0,1+\yshift) node[midway,scale=1.5] {$\vdots$};
%       \fi
%     \fi
%   }
  
%   % LABELS (All color removed)
%   \node[above=1,align=center,mygreen!60!black] at (N1-2.90) {input\\[-0.2em]layer};
%   \node[above=1,align=center,myblue!60!black] at (N3-2.90) {hidden layers};
%   \node[above=1,align=center,myred!60!black] at (N\Nnodlen-2.90) {output\\[-0.2em]layer};
  
% \end{tikzpicture}


